% Generated by scripts/assemble.py; corrections belong in source/corrections.json.
\examhead{61}{重庆统考 2026 年数学分析试卷}{60}
% p061-b12; source PDF page 61
\par\Needspace{5\baselineskip}
\blockanchor{p061-b12}
\majorquestion{一}计算题, 共 105 分.\par
% p061-b13; source PDF page 61
\blockanchor{p061-b13}
\question{1}求极限  $\lim_{x\to+\infty}\frac{\ln\sqrt{x}+\sin^{2}x}{x}$ .\par
% p061-b14; source PDF page 61
\blockanchor{p061-b14}
\question{2}设 $a > 0$, 求极限 $\lim_{n \to \infty} \sqrt{a} \sqrt[4]{a} \sqrt[8]{a} \cdots \sqrt[2^n]{a}$.\par
% p061-b15; source PDF page 61
\blockanchor{p061-b15}
\question{3}求极限 $\lim_{x\to +\infty}\left[x - x^3\ln \left(1 + \frac{1}{x^2}\right)\right]$\par
% p061-b16; source PDF page 61
\par\Needspace{6\baselineskip}
\blockanchor{p061-b16}
\question{4}求极限\par
\nopagebreak[4]
% p061-b17; source PDF page 61
\blockanchor{p061-b17}
\displayformula{\lim _ {n \rightarrow \infty} \left(\frac {1}{\sqrt {n ^ {2} + 1 ^ {2}}} + \frac {1}{\sqrt {n ^ {2} + 2 ^ {2}}} + \dots + \frac {1}{\sqrt {n ^ {2} + n ^ {2}}}\right).}
% p061-b18; source PDF page 61
\par\Needspace{6\baselineskip}
\blockanchor{p061-b18}
\question{5}已知函数\par
\nopagebreak[4]
% p061-b19; source PDF page 61
\blockanchor{p061-b19}
\displayformula{g (t) = \left\{ \begin{array}{l l} t + 2, & t \geq 0 \\ A t ^ {3} + B t ^ {2} + C t + D, & - r <   t <   0 \\ 0, & t \leq - r \end{array} \right.}
% p061-b20; source PDF page 61
\blockanchor{p061-b20}
连续可微, 其中 r > 0.\par
% p061-b21; source PDF page 61
\blockanchor{p061-b21}
(1) 求 A, B, C, D 的值.\par
% p061-b22; source PDF page 61
\blockanchor{p061-b22}
(2) 求导函数  $g'(t)$ .\par
% p061-b23; source PDF page 61
\blockanchor{p061-b23}
\question{6}已知函数 $f(x) = \int_{1}^{x^2} \frac{\sin t}{t} \, \mathrm{d}t$，求 $\int_{0}^{1} x(f(x) + 1) \, \mathrm{d}x$.\par
% p061-b24; source PDF page 61
\blockanchor{p061-b24}
\question{7}求幂级数 $\sum_{n=1}^{+\infty} \frac{x^n}{n+1}$ 的收敛域与和函数.\par
% p061-b25; source PDF page 61
\blockanchor{p061-b25}
\question{8}求由方程 $z \sin y + xe^z = 1$ 所确定的隐函数 $z = f(x, y)$ 在点 (1,0) 的全微分.\par
% p061-b26; source PDF page 61
\blockanchor{p061-b26}
\question{9}计算曲线积分 $\oint_{L} x e^{2xy} \, \mathrm{d}s$, 其中 $L$ 是以 $O(0,0), A(1,0), B(1,1)$ 为顶点的三角形.\par
% p062-b00; source PDF page 62
\par\Needspace{6\baselineskip}
\blockanchor{p062-b00}
\question{10}求曲面积分\par
\nopagebreak[4]
% p062-b01; source PDF page 62
\blockanchor{p062-b01}
\displayformula{\iint_ {S} x \mathrm{d} y \mathrm{d} z + y \mathrm{d} z \mathrm{d} x + z \mathrm{d} x \mathrm{d} y.}
% p062-b02; source PDF page 62
\blockanchor{p062-b02}
其中 $S$ 为曲面 $z = \sqrt{x^2 + y^2}$ 与 $x^{2} + y^{2} + z^{2} = 2$ 所围立体的表面，方向取外侧.\par
% p062-b03; source PDF page 62
\blockanchor{p062-b03}
\question{11}设 $F(\theta) = \int_{0}^{\pi} \ln(1 - \theta \cos x) \, \mathrm{d}x, \theta \in (-1, 1)$.\par
% p062-b04; source PDF page 62
\blockanchor{p062-b04}
(1) 求  $F'(\theta)$ .\par
% p062-b05; source PDF page 62
\blockanchor{p062-b05}
(2) 求 $F\left(\frac{\sqrt{3}}{2}\right)$ 的值.\par
% p062-b06; source PDF page 62
\par\Needspace{5\baselineskip}
\blockanchor{p062-b06}
\majorquestion{二}证明题, 共 45 分.\par
% p062-b07; source PDF page 62
\par\Needspace{6\baselineskip}
\blockanchor{p062-b07}
\question{12}点 $M_{1}(x_{1}, y_{1}), M_{2}(x_{2}, y_{2})$ 的距离定义为\par
\nopagebreak[4]
% p062-b08; source PDF page 62
\blockanchor{p062-b08}
\displayformula{\begin{array}{c} d _ {1} (M _ {1}, M _ {2}) = \sqrt {| x _ {2} - x _ {1} | ^ {2} + (y _ {2} - y _ {1}) ^ {2}}; \\ d _ {2} (M _ {1}, M _ {2}) = | x _ {2} - x _ {1} | + | y _ {2} - y _ {1} |. \end{array}}
% p062-b09; source PDF page 62
\blockanchor{p062-b09}
证明: 平面点列 $\{(x_{n}, y_{n})\}$ 按距离 $d_{1}$ 收敛到 $(x_{0}, y_{0})$ 的充要条件是 $\{(x_{n}, y_{n})\}$ 按距离 $d_{2}$ 收敛到 $(x_{0}, y_{0})$.\par
% p062-b10; source PDF page 62
\blockanchor{p062-b10}
\question{13}设 $a > 1$, 证明不等式: $2a^{\frac{1}{2}} \leq a^{x} + a^{1 - x} \leq 1 + a, x \in [0,1]$.\par
% p062-b11; source PDF page 62
\par\Needspace{6\baselineskip}
\blockanchor{p062-b11}
\question{14}已知\par
\nopagebreak[4]
% p062-b12; source PDF page 62
\blockanchor{p062-b12}
\displayformula{f (x, y) = \left\{ \begin{array}{l l} (x ^ {2} + y ^ {2}) \cos \frac {1}{x ^ {2} + y ^ {2}}, & x ^ {2} + y ^ {2} \neq 0; \\ 0, & x ^ {2} + y ^ {2} = 0. \end{array} \right.}
% p062-b13; source PDF page 62
\blockanchor{p062-b13}
(1) 证明: 函数 $f(x, y)$ 在 $(0, 0)$ 处可微.\par
% p062-b14; source PDF page 62
\blockanchor{p062-b14}
(2) 证明: 函数 $f(x, y)$ 的偏导数在 $(0, 0)$ 处不连续.\par
% p062-b15; source PDF page 62
\blockanchor{p062-b15}
\question{15}已知 $f(x)$ 在 $(a, +\infty)$ 上可导，且 $\lim_{x \to +\infty} f'(x) = A$.\par
% p062-b16; source PDF page 62
\blockanchor{p062-b16}
(1) 若 $A > 0$, 证明: $\lim_{x \to +\infty} f(x) = +\infty$.\par
% p062-b17; source PDF page 62
\blockanchor{p062-b17}
(2) 若 $A = 0$, 证明: $\lim_{x \to +\infty} \frac{f(x)}{x} = 0$.\par
